\documentclass[10pt, titlepage]{article}

\usepackage[margin=1in]{geometry}
\usepackage{fancyhdr}
\usepackage{amsfonts}
\usepackage{amsmath}
\usepackage{amssymb}
\usepackage{booktabs}
\usepackage{tikz}
\usepackage{pgfplots}
\usepackage{float}
\usepackage{bm}


\title{\normalsize PLUS Session Worksheet \\[1em] \Large \textit{MATH 1554:} Linear Algebra}
\author{Covering lessons: 7.3 \& 7.4}
\date{Fall 2023}

% ------- Set up header -------
\renewcommand{\sectionmark}[1]{\markright{#1}} % from https://tex.stackexchange.com/a/317393/227665
\pagestyle{fancy}
\fancyhead[L]{Fall 2023}
%\fancyhead[C]{\textit{\nouppercase{\rightmark}}}
\fancyhead[C]{\textit{Constrained Opt. (7.3) \& SVD (7.4)}}
\fancyhead[R]{\textbf{\thepage}}
% -------------------------------------------------------

% ------- Set up Main Idea box environment -------
\newcounter{mainideas}
\setcounter{mainideas}{1}
\newcommand{\mainidea}{\textbf{Main Idea \#\themainideas \stepcounter{mainideas}}}
\setlength\fboxsep{1em}
%\newcommand{\midea}[2]{\begin{center}\fbox{\parbox{#1}{\centering \mainidea #2}} \end{center}} % Could be environment but I couldn't bother

% The way this works is that the savebox collects whatever was between the minipage and the end of it, then in the last line of the definition of the environment
% we regurgigate whatever was collected into an fbox which we center.
\newsavebox{\mybox}
\newenvironment{mideea}[1]
{%
\begin{lrbox}{\mybox}%
	\begin{minipage}{#1\textwidth}
		{\centering \mainidea \par}

}
{\end{minipage}\end{lrbox}\begin{center}\fbox{\usebox{\mybox}}\end{center}}
% -------------------------------------------------------

% ------- Set up activity environment -------
\newcounter{question}
\setcounter{question}{1}
\newenvironment{newquestion}
{
	\noindent\textbf{Question \thequestion}
}
{
\stepcounter{question}
\par
}
% -------------------------------------------------------


% ------- Set up math stuff -------
\renewcommand{\vec}[1]{\bm{#1}}
\DeclareMathOperator{\vspan}{Span}
\DeclareMathOperator{\row}{row}
\DeclareMathOperator{\col}{col}
\DeclareMathOperator{\colspace}{Col}
\DeclareMathOperator{\rowspace}{Row}
\DeclareMathOperator{\nullspace}{Nul}
\DeclareMathOperator{\rank}{rank}
% -------------------------------------------------------

% ------- Set up circle -------
\newcommand{\bubble}{\tikz[baseline=-0.625ex]\draw[radius=0.5em, fill=white] (0,0) circle;}
\newcommand{\dbubble}{\tikz[baseline=-0.625ex]\draw[radius=0.5em, fill=black] (0,0) circle;}
% -------------------------------------------------------

\begin{document}
\maketitle

\section{\centering T \& F}
	\begin{newquestion}
		For the following statements, answer \textit{true} if the statement is always true, and \textit{false} otherwise.
		\begin{table}[!h]
			\centering
			\begin{tabular}{c@{\hspace{2.5em}}cp{0.75\textwidth}}
				true & false & \\
				\specialrule{\lightrulewidth}{0.25em}{1.25em}
				\bubble & \bubble & The quadratic form defined by \(Q(\vec{x}) = \vec{x}^TBB^T \vec{x}\) is always positive semidefinite. \\[0.75em]
				\bubble & \bubble & If  \(A\) is a symmetric matrix, then the corresponding eigenvalues of \(A\) are equal to the corresponding singular values of \(A\). \\[0.75em]
				\bubble & \bubble & Any matrix has a singular value decomposition.\\[0.75em]
				\bubble & \bubble & If \(Q\) is a positive definite quadratic form with \(Q(\vec{x}) = \vec{x}^TA\vec{x}\) and \(A = A^T\), then the largest eigenvalue of \(A\) is the maximum value \(Q\) takes on in its domain.\\[0.5em]
				\bubble & \bubble & Given \(A \in \mathbb{R}^{m \times n}\),  the largest singular value of \(A\) is the largest value \(Q(\vec{x}) = \|A\vec{x}\|\) takes on under the constraint \(\|\vec{x}\| = 1\).\\[0.5em]
				\bubble & \bubble & For any nonzero matrix \(C \in \mathbb{R}^{m \times n}\), \(F(\vec{x}) = \|C\vec{x}\|^2\) defines a quadratic form.\\[0.5em]
				\bottomrule
			\end{tabular}
		\end{table}
	\end{newquestion}

\section{\centering Short Answer Questions}
	\begin{newquestion}
		Try out the following questions.
		\begin{enumerate}
			\item Given the quadratic form \(Q(\vec{x}) = x_1^2 + 4x_2^2 + x_3^2 + 2x_1x_3\), there exists a change of variables \(\vec{x} = P\vec{y}\)
				that yields a new quadratic form \(Q'\) such that \(Q'\) takes on the same values as \(Q\), with the additional property that no cross-terms
				exist in \(Q'(\vec{y})\). Find \(Q'(\vec{y})\).
				\begin{enumerate}
					\item What is \(\max_{\|\vec{x}\|  = 1} Q(\vec{x})\)?
					\item For what input(s) \(\vec{x}\) is \(\max_{\|\vec{x}\| = 1} Q(\vec{x})\) attained?
					\item Say the answer for the above question is a vector \(\vec{m}\). (That is, \(Q(\vec{m}) = \max_{\|\vec{x}\| = 1} Q(\vec{x})\)). What is the largest value \(Q\) could take on if we impose the 
						constraints \[\|\vec{x}\| = 1 \qquad \text{\emph{and}} \qquad \vec{x} \cdot \vec{m} = 0?\]
				\end{enumerate}
			 
			\item Find the SVD of the matrix \[A = \begin{bmatrix}1 & -1 \\ 0 & 0 \\ 1 & -1\end{bmatrix}.\]
				\begin{enumerate}
					\item Find an orthonormal basis for 
						\begin{enumerate}
							\item \(\colspace A\), 
							\item \(\rowspace A\), 
							\item \(\nullspace A\), and 
							\item \(\nullspace A^T\).
						\end{enumerate}
					\item What is \(\rank A\)? How does this tie to the number of non-zero singular values?
				\end{enumerate}
			\item For the above matrix \(A\), what is \(\max_{\|\vec{x}\| = 1} \|A\vec{x}\|\)? At what input(s) does this maximum occur?
		\end{enumerate} 
	\end{newquestion}

\include{solutions}
\end{document}
